% s08-sec5.tex —— §5 主项（p.46–52）
\section{主项}
\srcpage{46}

如今，我们这一策略最有说服力的证据是：当 $G$ 半单且单连通时，它可以分离出主贡献。事实上，对任何满足我们条件的群，区分 $G=G_{\mathrm{der}}$ 与一般情形的只有形式上的问题。结果表明，这个主项由 $\hat\theta(0)$ 给出。

我们假定 $G=G_{\mathrm{der}}$。此外 $G$ 几乎处处拟分裂，甚至处处拟分裂。于是导出子群 $G_{\mathrm{der}}(F_v)$ 就是 $G_{\mathrm{der}}(F_v)$ 本身。\yizhu{原文如此。按上下文，此句疑本意是「$G(F_v)$ 的导出子群就是 $G_{\mathrm{der}}(F_v)$ 本身」；因为已假定 $G=G_{\mathrm{der}}$，后一读法才有信息量。译文照录原文。}{依据：原书 p.46}（我们没有去找最好的参考文献，[St2] 是一种可能。） 由此，$G_{\mathrm{der}}(\mathbb A_F)$ 的唯一一维自守表示是平凡表示 $\pi_0$。按公式 (1.3)，对 $^LG$ 的表示 $\rho$ 与 (1.14) 的函数 $f$，主项于是为
\begin{equation}\tag{5.1}
\prod_i\zeta^S(s+i)\prod_{v\in S}\operatorname{tr}\pi_v(f_v),
\end{equation}
其中变量 $s$ 如第 1 节。整数 $i$ 由与主幂幺元相联系的 $\operatorname{SL}(2)$ 到 $^LG$ 的嵌入 $\varphi$ 给出。诸 $i$ 是
\[
z\mapsto\rho\cdot\varphi\begin{pmatrix}z&0\\0&z^{-1}\end{pmatrix}
\]
的权的一半。表达式 (5.1) 等于 $\operatorname{tr}(\pi_0(f))$，从而等于
\[
\int_{G(\mathbb A_F)}f(g)\,dg=\prod_v\int_{G(F_v)}f_v(g_v)\,dg_v.
\]
下一个公式由公式 (3.31) 推出：
\begin{equation}\tag{5.2}
\int_{G(F_v)}f_v(g_v)\,dg_v=\int_{A(F_v)}\theta_v(b_v;1)\,db_v,
\end{equation}
更一般地推出
\begin{equation}\tag{5.3}
\int_{G(F_v)}\frac{L_v(s,\sigma_{T/G})}{L_v(1,\sigma_{T/G})}f_v(g_v)\,dg_v
=\int\theta_v(b_v;s)\,db_v=\hat\theta_v(0;s).
\end{equation}
我们强调，这个积分中的函数通过 $T=T_{g_v}$——正则元 $g_v$ 的中心化子——依赖于 $b_v$ 或 $g_v$。把它写成
\[
f_v(g_v;s)=\frac{L_v(s,\sigma_{T/G})}{L_v(1,\sigma_{T/G})}f_v(g_v)
\]
是方便的。

\srcpage{47}
我们指出，商 $L(s,\sigma_{T/G})/L(1,\sigma_{T/G})$ 在 $1\le s\le1+\epsilon$ 上有界，因为它是因子 $(1-\alpha/q_v)/(1-\alpha/q_v^s)$ 的乘积，其中 $|\alpha|=1$，诸 $\alpha$ 属于一个有限的单位根集合。唯一稍嫌微妙的情形是 $\alpha=1$，但此时当 $s\ge1$ 有 $1-1/q_v\le1-1/q_v^s$。

下面的引理与公式 (4.6) 等价，但形式更讨几何学家喜欢。

\begin{lemma}
\setcounter{theorem}{3}
设 $v\in|F|$ 是 $G$ 具有约化约化的位。此时态射 $c:G\to A$ 延拓到 $O_v$。假定体积形式 $\omega_G$ 与 $\omega_A$ 延拓为其 $O_v$-模型上处处非零的形式。设 $f_v$ 是紧群 $G(O_v)$ 的特征函数。则对一切 $b_v\in A^{\mathrm{rs}}(O_v)$，有
\begin{equation}\tag{5.5}
\theta_v(b_v;1)=q_v^{-\dim(G)+\dim(T)}\frac{|G(\kappa_v)|}{|T(\kappa_v)|}
=L_v(1,\sigma_{T/G})\frac{|G_{\mathrm{der}}(\kappa_v)|}{q_v^{\dim(G_{\mathrm{der}})}}.
\end{equation}
这里 $T$ 是 $b_v$ 上方的一个截面 $g_v\in G^{\mathrm{rs}}(O_v)$ 的中心化子，是定义在 $O_v$ 上的环面。
\end{lemma}

要证明的结果如下。

\begin{proposition}
\setcounter{theorem}{5}
有等式
\[
\lim_{S'\to\infty,\ s\downarrow1}\hat\theta_{S'}(0,s)=\int_{G(\mathbb A)}f(g)\,dg.
\]
无穷乘积
\[
\prod_v\int_{G(F_v)}f(g_v)\,dg_v
\]
收敛。此外，对给定的 $v$，
\[
\lim_{s\to1}\hat\theta_v(0,s)=\int_{G(F_v)}f_v(g_v)\,dg_v.
\]
\end{proposition}

因此要证明的是：对给定的有限 $S''$，
\[
\lim_{S'\to\infty,\ s\downarrow1}
\Bigl\{\prod_{S'-S''}\hat\theta_v(0,s)-\prod_{S'-S''}\int_{G(F_v)}f_v(g_v)\,dg_v\Bigr\}=0.
\]
自然地选取 $S''$，使得一切可能的麻烦都被排除。

这里有两个乘积，其中第二个等于
\[
\prod_{S'-S''}\frac{|G_{\mathrm{der}}(\kappa_v)|}{q_v^{\dim G_{\mathrm{der}}}}
=\prod_{S'-S''}\bigl(1+O(q_v^{-2})\bigr),
\]

\srcpage{48}
从而绝对收敛。回忆我们假定了 $G=G_{\mathrm{der}}$。要完成证明，只需对第一个乘积中的因子给出类似的估计
\begin{equation}\tag{5.7}
\hat\theta_v(0,s)=1+O(q_v^{-3/2}).
\end{equation}
我们发现函数 $f_v(g_v;s)$ 一致有界，且支在紧集 $G(O_v)$ 上。不过它们相当复杂，在接近由 $\Delta$ 的平方 $D$ 给出的除子时变得非常不规则。由于这个函数不变，它也是 Steinberg–Hitchin 底上的函数，我们称之为判别式。对群 $\operatorname{SL}(2)$，若特征多项式是 $X^2-bX+1$、特征值是 $\alpha$ 与 $\alpha^{-1}$，则 $D=(\alpha^2-1)(\alpha^{-2}-1)=4-b^2$ 是直线上的除子，含两个单点——至少当特征不等于 $2$ 时如此！本文回避了这一情形。

(5.7) 的证明基于以下观察。首先，在 $G(O_v)$ 的紧开子集 $G^{\mathrm{tvl}}(O_v)$（下标 "tvl" 意为横截）上——它由约化 $\bar g_v$ 正则、但不必半单、且判别式 $D(g_v)$ 的赋值小于或等于 $1$ 的元素 $g_v$ 组成——值 $f_v(g_v;s)$ 只依赖于 $\bar g_v$。其次，这个值可以用 $\kappa_v$ 上群 $G$ 的一个 $\ell$-进层方便地表达。最后，这个开集在 $G(O_v)$ 中的补集测度太小，不足以妨碍估计 (5.7)。

尽管几何学家不会感到需要，为便于研究迹公式或自守形式的专家理解论证，我们随着论证的展开，在群 $\operatorname{SL}(2)$ 的框架内、以具体的方式回忆所用的性质。例如对 $\operatorname{SL}(2)$，开集 $G^{\mathrm{tvl}}(O_v)$ 由如下条件定义：一个特征值 $\alpha$ 满足 $|1-\alpha^2|_v\le q_v^{-1/2}$。因此，若剩余特征不等于 $2$，被排除的 $\alpha$——更确切地说，需要更精细论证的 $\alpha$——是那些满足 $|1\mp\alpha|_v=1,q_v^{-1/2}$ 的。所以，一如通常情形，对 $\operatorname{SL}(2)$，"横截"意味着 $g$ 在 $\mathfrak g_g\backslash\mathfrak g$ 中的特征值\yizhu{原文写作 $\mathfrak g_g\backslash\mathfrak g$，按上下文应理解为商空间 $\mathfrak g/\mathfrak g_g$，即模去中心化子后的空间；译文照原文记号排印。}{依据：原书 p.48}取到与 $1$ 足够远的值。至于一般情形，精确定义要复杂些。

\begin{lemma}
\setcounter{theorem}{7}
$G^{\mathrm{tvl}}(O_v)$ 在 $G(O_v)$ 中的补集测度为 $O(q_v^{-2})$。
\end{lemma}

\srcpage{49}
\begin{proof}
设 $G_{\mathrm{sing}}$ 是 $G$ 中正则元（不论是否半单）开集 $G_{\mathrm{reg}}$ 的补集，它是 Zariski 闭集。由 Steinberg（[St1]）可知 $G_{\mathrm{sing}}$ 是余维三的闭集。例如对 $G=\operatorname{SL}(2)$，它是 $\{\pm I\}$。

判别式函数 $D$ 在 $T/W$ 上、从而在 Steinberg–Hitchin 底 $A$ 上定义既约除子 $\operatorname{div}D$。记 $\operatorname{div}D_{\mathrm{sing}}$ 为 $\operatorname{div}D$ 中判别式以至少二阶消失、从而其微分为零的闭子概形。这个闭集一般是余维二的，故当 $G=\operatorname{SL}(2)$ 时为空，但也有例外。除这些例外之外，它在 $G_{\mathrm{reg}}$ 中的原像至少余维二，因为态射 $c:G_{\mathrm{reg}}\to A$ 是光滑的。由于 $G_{\mathrm{sing}}$ 余维三，并集 $c^{-1}(\operatorname{div}D_{\mathrm{sing}})\cup G_{\mathrm{sing}}$ 是 $G$ 的至少余维二的闭子概形。

最简单的例外情形是特征为 $2$ 时的群 $\operatorname{SL}(2)$，此时 $D=b^2$，因为 $\operatorname{div}D_{\mathrm{sing}}$ 含点 $b=0$。要避免这类问题，只需求特征大于 $2$。此时以下论证成立。不过我们强调，所寻求的理论理应无例外地成立。遗憾的是目前我们没有完整证明，但已把小特征的若干情形留待另一次讨论。

对每个 $\bar g_v\in G(\kappa_v)$，约化为 $\bar g_v$ 的元素 $g_v\in G(O_v)$ 之集的测度为 $q_v^{-\dim(G)}$。因此，在我们的假设下，约化落在 $c^{-1}(\operatorname{div}D_{\mathrm{sing}})\cup G_{\mathrm{sing}}$ 中的 $g_v\in G(O_v)$ 之集，当 $q_v$ 趋于无穷时测度为 $O(q_v^{-2})$。

记 $G'$ 为 $G$ 中 $c^{-1}(\operatorname{div}D_{\mathrm{sing}})\cup G_{\mathrm{sing}}$ 的补开集。设 $g_v\in G(O_v)$ 的约化 $\bar g_v\in G'(\kappa_v)$。若 $\bar g_v$ 不属于 $c^{-1}(\operatorname{div}D-\operatorname{div}D_{\mathrm{sing}})$，则 $\bar g_v$ 的判别式赋值为零，$g_v$ 属于 $G^{\mathrm{tvl}}$。若 $\bar g_v\in c^{-1}(\operatorname{div}D-\operatorname{div}D_{\mathrm{sing}})$，则 $D$ 的微分在 $\bar g_v$ 的切空间上诱导一个非零线性形式。由此，由给定 $g_v$ 变形得到的、约化为 $\bar g_v$ 的元素所占比例为 $1/q_v$，于是约化为 $\bar g_v$ 的 $g_v\in G_{\mathrm{reg}}(O_v)-G^{\mathrm{tvl}}(O_v)$ 之集测度为 $q_v^{-\dim(G)-1}$。此外，被判别式零化的 $G$ 中点构成 $G$ 的一个除子。由于除子 $c^{-1}(\operatorname{div}D-\operatorname{div}D_{\mathrm{sing}})$ 余维一，$G^{\mathrm{tvl}}(O_v)$ 在 $G'(\kappa_v)$ 中的补集当 $q_v$ 趋于无穷时测度也是 $O(q_v^{-2})$。引理得证。
\end{proof}

我们指出，按定义 $G^{\mathrm{tvl}}$ 含于 $G^{\mathrm{rs}}(F_v)$，从而含于原像 $c^{-1}(A^{\mathrm{rs}})(F_v)$。考虑 Grothendieck–Springer 同时消解 $\pi:\tilde G\to G$，它在 $G^{\mathrm{rs}}$ 上的限制是有限态射 $\pi^{\mathrm{rs}}:\tilde G^{\mathrm{rs}}\to G^{\mathrm{rs}}$。与 $\tilde G^{\mathrm{rs}}$ 的一点相联系的是一个极大环面及包含它的 Borel 子群。显然，包含给定极大环面的 Borel 子群之集与 Weyl 群有典范等同。因此，直像 $\pi^{\mathrm{rs}}_*\bar{\mathbb Q}_\ell$ 是带 $W$-作用的 $\ell$-进层。我们可以引入层
\begin{equation}\tag{5.9}
\bigl(\pi^{\mathrm{rs}}_*\bar{\mathbb Q}_\ell\otimes X^*(T)\bigr)^W,
\end{equation}
但我们需要一个更精细的对象。

这要用到 [N] 中的结果，它们是在剩余特征不整除 Weyl 群阶的假设下证明的；对这篇不完整的文章而言，我们接受这一假设。把 $g$ 映到其半单部分与幂幺部分之积中的半单部分，给出映射 $G\to T/W$。还有映射 $\tilde G\to T$。它们给出

\srcpage{50}
\[
\begin{array}{ccc}
\tilde G&\longrightarrow&T\\
\downarrow&&\downarrow\\
G&\longrightarrow&T/W
\end{array}
\]
这个图限制后给出笛卡尔图
\[
\begin{array}{ccc}
\tilde G_{\mathrm{reg}}&\longrightarrow&T\\
\downarrow&&\downarrow\\
G_{\mathrm{reg}}&\longrightarrow&T/W
\end{array}
\]
论证如下。态射 $G_{\mathrm{reg}}\to T/W$ 光滑，且基变换后态射 $G_{\mathrm{reg}}\times_{T/W}T\to T$ 也光滑。因此纤维积也光滑，从而正规。有态射从 $\tilde G_{\mathrm{reg}}$ 到纤维积，它一方面是有限的（因为已知 $\tilde G_{\mathrm{reg}}\to G_{\mathrm{reg}}$ 有限），另一方面是双有理的（因为在 $G^{\mathrm{rs}}$ 上方这个图显然是笛卡尔的）。正规性于是推出态射
\[
\tilde G_{\mathrm{reg}}\to G_{\mathrm{reg}}\times_{T/W}T
\]
是同构。由此，群 $W$ 作用在 $\tilde G_{\mathrm{reg}}$ 上；这个作用不能延拓到 $\tilde G$。

于是可以引入层
\[
L=\bigl(\pi^{\mathrm{reg}}_*\bar{\mathbb Q}_\ell\otimes X^*(T)\bigr)^W.
\]
$L$ 在开集 $G^{\mathrm{rs}}$ 上的限制 (5.9) 是一个局部系统，其单值性由几何平展 Galois 覆盖

\srcpage{51}
$\tilde G^{\mathrm{rs}}\to G^{\mathrm{rs}}$ 以及 $W$ 在 $X^*(T)\otimes\bar{\mathbb Q}_\ell$ 上的通常作用给出。我们已假定 $G$ 半单。因此 $W$ 在 $X^*(T)$ 上的表示不含平凡表示，$L$ 也不以常层为子商。

由于 $G$ 单连通，正则半单元的中心化子连通，中心化子概形 $I$ 在 $G$ 上对开集 $G^{\mathrm{rs}}$ 的限制是环面 $I\times_GG^{\mathrm{rs}}$。这个环面的特征在 $G^{\mathrm{rs}}$ 上构成局部系统，它在把系数换到 $\bar{\mathbb Q}_\ell$ 之后同构于 $L$。我们没有找到这一断言的参考文献，但其在 Lie 代数情形的类比由 Donagi 与 Gaitsgory 在 [DG] 中证明。[N] 的命题 2.4.7 也是一个方便的参考。单连通群情形的证明与 Lie 代数情形完全类似。

若 $g_v:\operatorname{Spec}(O_v)\to G$ 是横截于判别式除子的 trait（即离散赋值环的谱所给出的一维正则局部概形），则 $\operatorname{Spec}(O_v)$ 在 $\tilde G$ 中的原像是若干 trait 的并，也就是 $\operatorname{Spec}(O_v)$ 的有限正规覆盖。证明这一断言的论证可在 [N] 4.7.3 的证明中找到。下面的引理描述了 $L$ 在 $g_v$ 的约化 $\bar g_v$ 处的纤维，它由此推出。

\begin{lemma}
\setcounter{theorem}{9}
设 $g_v\in G^{\mathrm{tvl}}(O_v)$。记 $X^*(T_{g_v})$ 为环面 $T_{g_v}$ 的特征模。它定义在 $F_v$ 的一个有限代数扩张上，带有 $F_v$ 的 Galois 群的连续作用。于是惯性群不变部分 $(X^*(T_{g_v})\otimes\bar{\mathbb Q}_\ell)^{I_v}$ 作为带 Frobenius 作用的 $\bar{\mathbb Q}_\ell$-向量空间同构于纤维 $L_{\bar g_v}$。
\end{lemma}

借助这个引理，我们可以用支撑在 $G_{\mathrm{reg}}(O_v)$ 上的函数 $f'_v(g_v,s)$ 代替函数 $f_v(g_v,s)$，它对 $g_v\in G_{\mathrm{reg}}(O_v)$ 取值
\[
f'_v(g_v,s)=1-\operatorname{tr}(\mathrm{Fr}_v,L_{\bar g_v})q_v^{-1}
+\operatorname{tr}(\mathrm{Fr}_v,L_{\bar g_v})q_v^{-s}
=1-\operatorname{tr}(\mathrm{Fr}_v,L_{\bar g_v})\bigl(q_v^{-1}-q_v^{-s}\bigr).
\]
只需比较两个函数在开集 $G^{\mathrm{tvl}}(O_v)$ 上的积分，因为其中任一个在其补集上的积分都是 $O(q_v^{-2})$。在 $G^{\mathrm{tvl}}(O_v)$ 上，它们的差也被 $q_v^{-2}$ 界住。

为证明命题 5.6，我们首先注意到
\[
\int_{G^{\mathrm{tvl}}(O_v)}dg_v=\int_{G(O_v)}dg_v+O(q_v^{-2})
=q_v^{-\dim G}|G(\kappa_v)|+O(q_v^{-2}),
\]
且
\begin{equation}\tag{5.11}
q_v^{-\dim G}|G(\kappa_v)|=1+O(q_v^{-2}).
\end{equation}

\srcpage{52}
最后一个等式见于 Ono 关于 Tamagawa 数的文章。

由于 $G_{\mathrm{reg}}$ 的余维为 $3$，为完成命题 5.6 的证明，只需验证估计
\begin{equation}\tag{5.12}
q_v^{-\dim G}\sum_{\bar g_v\in G_{\mathrm{reg}}(\kappa_v)}\operatorname{tr}(\mathrm{Fr}_v,L_{\bar g_v})
=O(q_v^{-1/2}).
\end{equation}
这是下面引理的直接推论。

\begin{lemma}
\setcounter{theorem}{12}
当 $q_v\to\infty$ 时，有估计
\[
\sum_{\bar g_v\in G_{\mathrm{reg}}(\kappa_v)}\operatorname{tr}(\mathrm{Fr}_v,L_{\bar g_v})
=O\bigl(q_v^{\dim(G)-1/2}\bigr).
\]
\end{lemma}

对 $G_{\mathrm{reg}}$ 限制到 $\operatorname{Spec}\kappa_v$ 上所得的概形上的层 $\mathcal L$，由 Grothendieck–Lefschetz 迹公式，上述表达式等于
\[
\sum_{i=0}^{2\dim(G)}(-1)^i\operatorname{tr}\bigl(\mathrm{Fr},H^i_c(G_{\mathrm{reg}}\otimes_{O_v}\bar\kappa_v,L)\bigr).
\]
注意，除有限个位外，这些上同调群的维数与位 $v$ 无关。按 [D] 的主要定理，Frobenius 在 $H^i_c(G_{\mathrm{reg}}\otimes_{O_v}\bar\kappa_v,L)$ 上的特征值的绝对值以 $q_v^{i/2}$ 为界。因此只需证明 $H^{2\dim(G)}_c(G_{\mathrm{reg}},L)=0$。但这来自 $L$ 不以平凡局部系统为子商这一事实。

对群 $\operatorname{SL}(2)$ 与奇剩余特征，命题 5.6 的证明可以用初等方法完成。$F_v$ 上有三类环面：分裂、非分歧与分歧。它们对常函数 $1$ 的积分的贡献形如
\[
a_1+a_2q_v^{-1}+O(q_v^{-2}),\qquad
b_1+b_2q_v^{-1}+O(q_v^{-2}),\qquad
c_2q_v^{-1}+O(q_v^{-2}).
\]
可以手算验证 $a_1+b_1=1$ 与 $a_2+b_2+c_2=0$。它们对 $q_v^{-1}-q_v^{-s}$ 的系数的贡献分别形如 $a_3$、$b_3$ 与 $0$。再手算验证 $a_3+b_3=O(q_v^{-1})$。一般情形下，正是这第二个计算需要精细得多的处理——Grothendieck–Lefschetz 公式。
